\section{Neural Compute Unit, memoria de largo plazo y personalización}

Ver el LLM como una neural compute unit es la analogía de sistemas más importante de la segunda mitad de la clase. El modelo, como una unidad de cómputo, recibe token instructions y produce tokens, pero no obtiene de forma automática persistent memory, process isolation, device driver ni permission model. Solo al entender el alcance de la analogía se puede diseñar correctamente una agent architecture.

\subsection{Model as compute unit}

En la clase se comparan los tokens de lenguaje natural con las binary instructions. Lo que tienen en común es que la entrada pasa por un cómputo y se convierte en salida; la diferencia es que el comportamiento de un LLM es probabilístico, su espacio semántico es abierto y su entrada y salida carecen del tipado estricto y del determinismo de una ISA tradicional. Al leer la figura conviene mantener a la vez la compute analogy y las missing systems guarantees.

\lecturefigure{slide-090.jpg}{El AI model como neural compute unit token-in/token-out}{Presentación oficial de CS25 V4, página 90}

\begin{knowledgebox}{Lectura de la figura: límites de validez de la CPU analogy}
\begin{tabularx}{\textwidth}{L{0.22\textwidth} L{0.32\textwidth} >{\raggedright\arraybackslash}X}
\toprule
Nivel & Dónde la analogía es válida & Dónde la analogía falla \\
\midrule
Instruction & El prompt/context especifica el cómputo de esta ronda & El lenguaje natural es ambiguo y carece de opcodes fijos y de sistema de tipos. \\
Compute & Los parámetros ejecutan operaciones matriciales a gran escala & La salida es una distribución de probabilidad; no se pueden suponer garantías de determinismo ni de corrección. \\
Memory & El context ofrece un espacio de trabajo & El context tiene un límite de longitud y no se vuelve persistente de forma automática. \\
I/O & Las tools amplían las capacidades externas & Los permisos, las excepciones y los efectos secundarios deben gestionarse explícitamente en el sistema. \\
\bottomrule
\end{tabularx}
\end{knowledgebox}

\teachervoice{El ponente llama al modelo neural compute unit e imagina escribir un nuevo programming language para el Transformer. El valor didáctico de esta analogía está en convertir el prompt, la memory, las tools y las variables en componentes de un sistema, en lugar de atribuirle al modelo rasgos humanos como si fuera un cerebro completo.}

\subsection{Looped Transformer y cómputo repetido}

Un forward pass corresponde a un cómputo de profundidad fija. La diapositiva del Looped Transformer propone bloques compartidos o repetidos para que el modelo ejecute más iteraciones con los mismos parámetros. Se parece al reasoning loop: ambos intercambian más computation steps por un procesamiento más complejo, pero deben resolver la termination y la stability.

\lecturefigure{slide-091.jpg}{Looped Transformer y módulos de cómputo repetido}{Presentación oficial de CS25 V4, página 91}

\begin{importantbox}{Iterative computation}
Si el mismo bloque $F_\theta$ se repite $K$ veces:
\[
h^{(k+1)}=F_\theta(h^{(k)},x),\qquad k=0,\ldots,K-1.
\]
Compartir parámetros reduce la velocidad de crecimiento del modelo, pero puede traer problemas de optimization, de fixed-point stability y de stopping dinámico. Repetir más capas no equivale automáticamente a un reasoning más profundo.
\end{importantbox}

\subsection{La long-term memory no es solo vector search}

En la clase se compara la persistent memory con un disco: primero se codifican los documentos o las experiencias como embeddings y se almacenan, y cuando hace falta se recuperan hacia el context. Los sistemas actuales suelen usar nearest-neighbor search, pero una memoria de largo plazo real también requiere tiempo, jerarquía, relaciones y estrategias de actualización. Esta sección explica primero el embedding pipeline actual y después amplía el memory schema con las open questions planteadas en la clase.

\lecturefigure{slide-092.jpg}{Problemas de almacenamiento, embedding y recuperación en la long-term memory}{Presentación oficial de CS25 V4, página 92}

\begin{importantbox}{Memory retrieval score}
Para un query embedding $q$ y un memory item $m_i$, la similitud más simple puede escribirse como
\[
s_i=\frac{q^\top m_i}{\lVert q\rVert_2\lVert m_i\rVert_2},\qquad
\mathcal{I}=\operatorname{TopK}_i(s_i).
\]
Esto solo aprovecha la similitud semántica. Un memory ranking real también debería incorporar recency, importance, user scope, causal link, task state y access control.
\end{importantbox}

\begin{knowledgebox}{Concepto previo: los cuatro niveles de la agent memory}
\begin{enumerate}
  \item \textbf{Working memory}: el context actual, el scratchpad y los tool results;
  \item \textbf{Episodic memory}: interactions y action trajectories pasadas;
  \item \textbf{Semantic memory}: hechos estables, documentos y conocimiento del usuario;
  \item \textbf{Procedural memory}: workflows, policies y skills reutilizables.
\end{enumerate}
Construir una sola base de datos vectorial normalmente no basta para cubrir a la vez estas cuatro semánticas.
\end{knowledgebox}

\teachervoice{El ponente señala que la memory actual suele quedarse en un k-NN simple. La coherencia temporal, la graph structure, la online adaptation y los datos en cambio constante no quedan resueltos de forma natural; «conectar una vector database» es solo el punto de partida.}

\subsection{Personalization: las preferencias cambian y los datos tienen límites}

La diapositiva de Personalization aspira a que el agent sea una extensión digital del usuario. Para ello debe aprender su estilo de escritura, su agenda, su tolerancia al riesgo y sus objetivos de largo plazo, y al mismo tiempo distinguir entre preferencias estables, instrucciones de una sola vez e información sensible. Las dos diapositivas siguientes deben leerse considerando a la vez la capability y el data boundary, con especial atención a cómo se recopilan, actualizan, revocan y sobrescriben las preferencias.

\lecturefigure{slide-093.jpg}{Objetivos de la user-agent alignment y la personalization}{Presentación oficial de CS25 V4, página 93}

\lecturefigure{slide-094.jpg}{Retos de recopilación de preferencias, privacidad, cambio y alignment}{Presentación oficial de CS25 V4, página 94}

\teachervoice{El ponente recurre al hippocampus y a los datos personales en cambio constante para recordar al lector que la memory no es un repositorio estático de documentos. El sistema debe manejar el orden temporal, la estructura de relaciones, la actualización en línea y los cambios en las preferencias del usuario.}

\begin{warningbox}{Cuatro fallas comunes de la Personalization}
\begin{enumerate}
  \item Interpretar erróneamente una conducta aislada como una preferencia permanente;
  \item Mezclar la memory de distintos usuarios o workspaces;
  \item Guardar datos sensibles sin límite para «entender mejor al usuario»;
  \item Obedecer ciegamente cuando las preferencias del usuario entran en conflicto con la seguridad, la ley o la policy de la organización.
\end{enumerate}
El sistema necesita consent, scope, retention, edit/delete y policy precedence.
\end{warningbox}

\begin{importantbox}{La preference update requiere decaimiento temporal y confirmación}
El score de una preference $p$ puede escribirse como
\[
S(p)=\sum_j w_j\,e^{-\lambda\Delta t_j}\,\operatorname{evidence}_j(p),
\]
donde $w_j$ indica la confiabilidad de la evidencia, $\Delta t_j$ es el intervalo de tiempo y $\lambda$ controla el decaimiento. Las preferencias de alto riesgo no deben inferirse solo a partir de conductas implícitas; deben requerir confirmación explícita.
\end{importantbox}

\subsection{Resumen de la sección}

La neural-compute analogy aclara los límites entre model, context, memory y tool, pero no permite tratar al LLM como una CPU con una ISA determinista y garantías de aislamiento. La memoria de largo plazo necesita jerarquía, tiempo y permisos, y la personalización necesita mecanismos de consent, actualización y borrado. Cuanto más cerca del usuario están las capacidades del sistema, más se convierte la memory governance en un elemento central del diseño.
